\chapter{Introduction: FreeCAD's Architecture}
\label{ch:introduction}

FreeCAD is a free, open-source, general-purpose parametric 3D CAD
modeler. For the engineer coming from a commercial package
(SolidWorks, Inventor, CATIA), three architectural ideas explain almost
everything about how FreeCAD behaves and why. This short chapter fixes
them before we touch a mouse.

\section{Parametric, history-based modeling}

A FreeCAD model is not a static mesh of triangles; it is a
\emph{recipe}. Each operation --- a sketch, a pad, a fillet, a Boolean
cut --- is stored as a \emph{feature} that remembers its parameters and
its inputs. Re-evaluating the recipe (\emph{recomputing} the document)
regenerates the geometry. Change a sketch dimension and every downstream
feature updates. This is the same \emph{feature tree} paradigm used by
mainstream MCAD, and it is what makes CAD models editable rather than
merely drawable.

The underlying geometry kernel is \textbf{OpenCascade Technology (OCCT)},
which provides the boundary-representation (B-rep) solids, the NURBS
surfaces, and the geometric operations. FreeCAD is, in large part, an
application framework wrapped around OCCT plus a scripting layer.

\section{The document object model}

Everything in a FreeCAD file (\texttt{.FCStd}) lives in a
\textbf{Document}, which is a container of \textbf{Document Objects}. A
document object (\texttt{App::DocumentObject}) is a typed node --- a
sketch, a solid, a mesh, an FEM constraint --- with:

\begin{itemize}
  \item a set of \textbf{Properties} (typed, introspectable fields:
        lengths, placements, links to other objects, \ldots);
  \item a \textbf{dependency graph} linking it to the objects it consumes;
  \item an \texttt{execute()} method that recomputes it from its inputs.
\end{itemize}

The objects and their links form a \emph{directed acyclic graph}; a
recompute is a topological traversal of that graph. The \emph{Model} tab
of the tree view is simply a human-readable rendering of this object
graph.

\begin{fctip}[Scripting is first-class]
Every property you set through the GUI is equally reachable from Python.
The document object model \emph{is} the API:
\begin{lstlisting}[language=FreeCADPython]
import FreeCAD as App
doc = App.newDocument("Demo")
box = doc.addObject("Part::Box", "Box")
box.Length = 40      # a property, identical to the GUI field
doc.recompute()
\end{lstlisting}
This matters for FEA: complete analyses can be built, meshed, solved and
post-processed headlessly in Python, which is how automated verification
and parameter sweeps are done (Part~III).
\end{fctip}

\section{The workbench paradigm}

FreeCAD does not present one monolithic toolset. Instead it groups
functionality into \textbf{Workbenches} --- self-contained sets of tools,
toolbars and menus for a particular task. \textsf{Sketcher} constrains 2D
profiles; \textsf{PartDesign} builds solids feature-by-feature;
\textsf{Part} offers direct primitive-and-Boolean modeling;
\textsf{FEM} performs finite-element analysis; \textsf{TechDraw}
produces drawings, and so on. Switching workbench changes the available
tools but \emph{not} the document --- all workbenches operate on the same
underlying object graph.

For this book the relevant workbenches are \textsf{Sketcher},
\textsf{Part}, \textsf{PartDesign} (Part~I) and \textsf{FEM}
(Part~III).

\section{Scope of this book}

We assume FreeCAD~1.x. Part~I is deliberately brisk: enough of the
interface, Sketcher and solid modeling to build the geometry an analysis
needs, and no more. Readers wanting exhaustive modeling tutorials should
consult the official wiki~\citep{freecad_wiki}; readers who already model
in FreeCAD can jump to the FEM theory in Part~II.
